\section{Mô hình PhoWhisper}

PhoWhisper là bộ mô hình nhận dạng giọng nói tiếng Việt được phát triển
bởi VinAI Research, dựa trên kiến trúc OpenAI Whisper và được fine-tune
trên dữ liệu tiếng Việt quy mô lớn. Mô hình hỗ trợ nhận dạng chính xác
các đặc thù ngữ âm tiếng Việt bao gồm 6 thanh điệu, các âm đặc trưng
và sự khác biệt giữa giọng vùng miền (Bắc, Trung, Nam).

Dự án sử dụng phiên bản PhoWhisper-large (\texttt{vinai/PhoWhisper-large}),
mô hình lớn nhất trong bộ PhoWhisper với kích thước khoảng 2 GB. Model
được triển khai local trên server thông qua HuggingFace Transformers
pipeline \texttt{automatic-speech-recognition}, không phụ thuộc cloud API.
So với phiên bản nhỏ hơn (PhoWhisper-tiny), phiên bản large cho kết quả
nhận dạng chính xác hơn đáng kể, đặc biệt với giọng vùng miền và môi
trường có nhiễu.

Trong quá trình phát triển, module STT đã trải qua nhiều lần nâng cấp:
PhoWhisper-tiny (local, nhẹ nhưng accuracy thấp) $\rightarrow$ whisper-1
(OpenAI cloud, accuracy tốt nhưng phụ thuộc internet) $\rightarrow$
PhoWhisper-large (local, accuracy cao và không phụ thuộc cloud).

\section{Mô hình GPT-4o-mini}

GPT-4o-mini là mô hình ngôn ngữ lớn (LLM) nhỏ gọn và tiết kiệm chi phí
của OpenAI, được thiết kế cho các tác vụ không yêu cầu mô hình đầy đủ
GPT-4o. Trong dự án, GPT-4o-mini được sử dụng để sinh mô tả sản phẩm
từ text nhận dạng giọng nói thông qua OpenAI API.

Mô hình có khả năng hiểu ngữ cảnh sâu, cho phép xử lý hiệu quả text
lộn xộn từ STT --- nhận diện từ lấp (``ừm'', ``à''), bỏ qua các đoạn
lặp lại và lỗi phát âm, trích xuất thông tin cốt lõi về sản phẩm và
sinh văn bản mạch lạc theo phong cách viết được chỉ định. Với temperature
0.75 và giới hạn 600 tokens, mô hình cân bằng giữa tính sáng tạo (văn
phong đa dạng) và tính nhất quán (nội dung chính xác).
